% GENERATED FILE. Edit canonical lessons, then run npm run generate:books.
% canonical-source-hash: fnv1a64:61db040ed69c6ded
% canonical-lessons: AR-C318-hadara, AR-C318-athar, AR-C318-atlal, AR-C318-mabad, AR-C318-haram

\chapter{Ancient civilizations}
\label{ch:ar-ancient-civilizations}
\hlchaptermodality{318}

\begin{chapteropening}
\textbf{By the end of this chapter:} \emph{I can say a civilization, a trace, remains, a temple and a pyramid.}

Complete the last lesson of chapter 318: I can say a civilization, a trace, remains, a temple and a pyramid.
\end{chapteropening}

\section[\texorpdfstring{\ar{حضارة} (ḥaḍāra)}{حضارة (ḥaḍāra)}]{\ar{حضارة} (ḥaḍāra) — a civilization}

\label{lesson:AR-C318-hadara}



\begin{quote}
Before the new one: say the Arabic for a patrol, then the Arabic for a checkpoint.
\end{quote}

\subsection*{You'll want to know: \ar{حضارة}}
\textbf{\ar{حضارة}} — \emph{ḥaḍāra} — \textquotedblleft{}a civilization\textquotedblright{}.

\begin{quote}
\ar{درسنا} \ar{حضارة} \ar{قديمة} \ar{في} \ar{المدرسة}.
\end{quote}

\begin{quote}
\emph{darasnā ḥaḍāra qadīma fī al-madrasa.}
\end{quote}

\begin{quote}
We studied an ancient civilization at school.
\end{quote}

\begin{cousinweb}
\textbf{\ar{حضارة}} (\emph{ḥaḍāra}) is built on the native Arabic root \emph{ḥ-ḍ-r} (Lane), \textquotedblleft{}to be present\textquotedblright{}, the root of \emph{ḥaḍara}, to attend, and of the taught \emph{ḥaḍḍara}, to prepare. Lane gives \emph{ḥaḍāra} as dwelling in towns and villages, the opposite of desert life.
\end{cousinweb}

\begin{grammarlens}[title={What you've built}]
One more word for telling what happened.
\end{grammarlens}

\subsection*{Your turn}
\begin{itemize}
\raggedright
  \item \emph{Say it:} \emph{ḥaḍāra}
  \item \emph{Say it:} \emph{ḥaḍāra}, once more
  \item \emph{Say it:} \emph{darasnā ḥaḍāra qadīma fī al-madrasa.} — We studied an ancient civilization at school.
  \item \emph{Recall:} say the Arabic for a patrol, then the Arabic for a checkpoint, then say \emph{ḥaḍāra} again
\end{itemize}

\begin{tcolorbox}[breakable,colback=teal!4,colframe=teal!35!black,title={Before you move on}]
What does \ar{حضارة} mean? (\textquotedblleft{}A civilization\textquotedblright{}.) Say it once more.
\end{tcolorbox}
\section[\texorpdfstring{\ar{أثر} (athar)}{أثر (athar)}]{\ar{أثر} (athar) — a trace, a mark; an effect}

\label{lesson:AR-C318-athar}



\begin{quote}
Before the new one: say the Arabic for a checkpoint, then the Arabic for a civilization.
\end{quote}

\subsection*{You'll want to know: \ar{أثر}}
\textbf{\ar{أثر}} — \emph{athar} — \textquotedblleft{}a trace, a mark; an effect\textquotedblright{}.

\begin{quote}
\ar{رأينا} \ar{أثر} \ar{قدم} \ar{في} \ar{الرمل}.
\end{quote}

\begin{quote}
\emph{raʾaynā athar qadam fī ar-raml.}
\end{quote}

\begin{quote}
We saw a footprint in the sand.
\end{quote}

Plural: \ar{آثار}, \emph{āthār}.

The plural \ar{الآثار} is antiquities: \ar{عالم} \ar{آثار}, an archaeologist.

\begin{cousinweb}
\textbf{\ar{أثر}} (\emph{athar}) is built on the native Arabic root \emph{ʾ-th-r} (Lane), where \emph{athar} is a mark or track left behind; \emph{athara} is also \textquotedblleft{}to hand down a report\textquotedblright{}.
\end{cousinweb}

\begin{grammarlens}[title={What you've built}]
One more word for telling what happened.
\end{grammarlens}

\subsection*{Your turn}
\begin{itemize}
\raggedright
  \item \emph{Say it:} \emph{athar}
  \item \emph{Say it:} \emph{athar}, once more
  \item \emph{Say it:} \emph{raʾaynā athar qadam fī ar-raml.} — We saw a footprint in the sand.
  \item \emph{Recall:} say the Arabic for a checkpoint, then the Arabic for a civilization, then say \emph{athar} again
\end{itemize}

\begin{tcolorbox}[breakable,colback=teal!4,colframe=teal!35!black,title={Before you move on}]
What does \ar{أثر} mean? (\textquotedblleft{}A trace, a mark; an effect\textquotedblright{}.) Say it once more.
\end{tcolorbox}
\section[\texorpdfstring{\ar{أطلال} (aṭlāl)}{أطلال (aṭlāl)}]{\ar{أطلال} (aṭlāl) — remains of a dwelling, ruins}

\label{lesson:AR-C318-atlal}



\begin{quote}
Before the new one: say the Arabic for a civilization, then the Arabic for a trace.
\end{quote}

\subsection*{You'll want to know: \ar{أطلال}}
\textbf{\ar{أطلال}} — \emph{aṭlāl} — \textquotedblleft{}remains of a dwelling, ruins\textquotedblright{}.

\begin{quote}
\ar{وقف} \ar{الشاعر} \ar{على} \ar{أطلال} \ar{البيت}.
\end{quote}

\begin{quote}
\emph{waqafa ash-shāʿir ʿalā aṭlāl al-bayt.}
\end{quote}

\begin{quote}
The poet stood at the remains of the house.
\end{quote}

The old Arabic ode opens with the poet weeping at the \ar{أطلال} of a camp.

\begin{cousinweb}
\textbf{\ar{أطلال}} (\emph{aṭlāl}) is built on the native Arabic root \emph{ṭ-l-l} (Lane), where \emph{ṭalal} is the part of a dwelling that remains standing after people leave; \emph{aṭlāl} is its plural.
\end{cousinweb}

\begin{grammarlens}[title={What you've built}]
One more word for telling what happened.
\end{grammarlens}

\subsection*{Your turn}
\begin{itemize}
\raggedright
  \item \emph{Say it:} \emph{aṭlāl}
  \item \emph{Say it:} \emph{aṭlāl}, once more
  \item \emph{Say it:} \emph{waqafa ash-shāʿir ʿalā aṭlāl al-bayt.} — The poet stood at the remains of the house.
  \item \emph{Recall:} say the Arabic for a civilization, then the Arabic for a trace, then say \emph{aṭlāl} again
\end{itemize}

\begin{tcolorbox}[breakable,colback=teal!4,colframe=teal!35!black,title={Before you move on}]
What does \ar{أطلال} mean? (\textquotedblleft{}Remains of a dwelling, ruins\textquotedblright{}.) Say it once more.
\end{tcolorbox}
\section[\texorpdfstring{\ar{معبد} (maʿbad)}{معبد (maʿbad)}]{\ar{معبد} (maʿbad) — a temple}

\label{lesson:AR-C318-mabad}



\begin{quote}
Before the new one: say the Arabic for a trace, then the Arabic for remains of a dwelling.
\end{quote}

\subsection*{You'll want to know: \ar{معبد}}
\textbf{\ar{معبد}} — \emph{maʿbad} — \textquotedblleft{}a temple\textquotedblright{}.

\begin{quote}
\ar{زرنا} \ar{معبدا} \ar{قديما} \ar{على} \ar{النهر}.
\end{quote}

\begin{quote}
\emph{zurnā maʿbadan qadīman ʿalā an-nahr.}
\end{quote}

\begin{quote}
We visited an ancient temple on the river.
\end{quote}

Plural: \ar{معابد}, \emph{maʿābid}.

\begin{cousinweb}
\textbf{\ar{معبد}} (\emph{maʿbad}) is built on the native Arabic root \emph{ʿ-b-d} (Lane), \textquotedblleft{}to serve, to worship\textquotedblright{}. The \emph{mafʿal} pattern names a place: \textquotedblleft{}a place of worship\textquotedblright{}.
\end{cousinweb}

\begin{grammarlens}[title={What you've built}]
One more word for telling what happened.
\end{grammarlens}

\subsection*{Your turn}
\begin{itemize}
\raggedright
  \item \emph{Say it:} \emph{maʿbad}
  \item \emph{Say it:} \emph{maʿbad}, once more
  \item \emph{Say it:} \emph{zurnā maʿbadan qadīman ʿalā an-nahr.} — We visited an ancient temple on the river.
  \item \emph{Recall:} say the Arabic for a trace, then the Arabic for remains of a dwelling, then say \emph{maʿbad} again
\end{itemize}

\begin{tcolorbox}[breakable,colback=teal!4,colframe=teal!35!black,title={Before you move on}]
What does \ar{معبد} mean? (\textquotedblleft{}A temple\textquotedblright{}.) Say it once more.
\end{tcolorbox}
\section[\texorpdfstring{\ar{هرم} (haram)}{هرم (haram)}]{\ar{هرم} (haram) — a pyramid}

\label{lesson:AR-C318-haram}



\begin{quote}
Before the new one: say the Arabic for remains of a dwelling, then the Arabic for a temple.
\end{quote}

\subsection*{You'll want to know: \ar{هرم}}
\textbf{\ar{هرم}} — \emph{haram} — \textquotedblleft{}a pyramid\textquotedblright{}.

\begin{quote}
\ar{الهرم} \ar{الكبير} \ar{قديم} \ar{جدا}.
\end{quote}

\begin{quote}
\emph{al-haram al-kabīr qadīm jiddan.}
\end{quote}

\begin{quote}
The Great Pyramid is very old.
\end{quote}

Plural: \ar{أهرام}, \emph{ahrām}.

\begin{cousinweb}
\textbf{\ar{هرم}} (\emph{haram}) shares its letters with the root \emph{h-r-m}, \textquotedblleft{}to grow very old\textquotedblright{} (Lane). Whether the pyramid's name comes from it or from another language is not settled.
\end{cousinweb}

\begin{grammarlens}[title={What you've built}]
One more word for telling what happened.
\end{grammarlens}

\subsection*{Your turn}
\begin{itemize}
\raggedright
  \item \emph{Say it:} \emph{haram}
  \item \emph{Say it:} \emph{haram}, once more
  \item \emph{Say it:} \emph{al-haram al-kabīr qadīm jiddan.} — The Great Pyramid is very old.
  \item \emph{Recall:} say the Arabic for remains of a dwelling, then the Arabic for a temple, then say \emph{haram} again
\end{itemize}

\begin{tcolorbox}[breakable,colback=teal!4,colframe=teal!35!black,title={Before you move on}]
What does \ar{هرم} mean? (\textquotedblleft{}A pyramid\textquotedblright{}.) Say it once more.
\end{tcolorbox}
